% !TEX program = xelatex
% 编译: xelatex 问题一_严格数学推导.tex
\documentclass[11pt,a4paper]{ctexart}

\usepackage{amsmath,amssymb,amsthm}
\usepackage{geometry}
\usepackage{booktabs}
\usepackage{array}
\usepackage{enumitem}
\usepackage[hidelinks]{hyperref}

\geometry{left=2.2cm,right=2.2cm,top=2.4cm,bottom=2.4cm}

\setlength{\parindent}{0pt}
\setlength{\parskip}{6pt plus 1pt}

\newtheorem{theorem}{定理}
\newtheorem{lemma}[theorem]{引理}
\newtheorem{proposition}[theorem]{命题}
\newtheorem{corollary}[theorem]{推论}
\theoremstyle{definition}
\newtheorem{definition}[theorem]{定义}
\theoremstyle{remark}
\newtheorem{remark}[theorem]{注记}

\newcommand{\RR}{\mathbb{R}}
\newcommand{\argvec}{\arg}
\newcommand{\diam}{\operatorname{diam}}
\newcommand{\wrap}{\operatorname{wrap}}
\newcommand{\MEC}{\operatorname{MEC}}

\title{\bfseries 问题一：交会定位区域及其直径的严格数学推导}
\author{}
\date{}

\begin{document}
\maketitle
\vspace{-2.5em}

\section{记号与前提}

\textbf{给定}：检测点 $S_1,\dots,S_n\in\RR^2$ 及相应的示向度 $\theta_1,\dots,\theta_n\in[0^\circ,360^\circ)$；测向误差界 $\tau=1^\circ$。

\textbf{前提（题设）}：设干扰源真实位置为 $G$，真实方位角 $\beta_i=\argvec(G-S_i)$，则
\begin{equation}
\bigl|\wrap(\beta_i-\theta_i)\bigr|\le\tau,\qquad
\wrap(x)=x-360^\circ\Bigl\lfloor\frac{x+180^\circ}{360^\circ}\Bigr\rfloor\in(-180^\circ,180^\circ].
\tag{A1}
\end{equation}

\section{定位区域及其凸性}

\begin{definition}
单个检测点确定的\textbf{视界楔形}为
\begin{equation}
W_i=\Bigl\{P:\ P=S_i\ \text{或}\ \bigl|\wrap\bigl(\argvec(P-S_i)-\theta_i\bigr)\bigr|\le\tau\Bigr\},
\end{equation}
\textbf{交会定位区域}为
\begin{equation}
\mathcal R=\bigcap_{i=1}^{n}W_i .
\end{equation}
\end{definition}

\begin{proposition}[相容性]
$G\in\mathcal R$。故 $\mathcal R$ 是干扰源位置的可行域。
\end{proposition}

\begin{proof}
由 (A1) 逐点得 $G\in W_i$，故 $G\in\bigcap_i W_i=\mathcal R$。
\end{proof}

\begin{lemma}[楔形的纵向--横向刻画]\label{lem:cone}
取标准正交基
\begin{equation}
u_i=(\cos\theta_i,\ \sin\theta_i),\qquad v_i=(-\sin\theta_i,\ \cos\theta_i),
\end{equation}
记 $d=P-S_i$，$\xi_i=d\cdot u_i$，$\eta_i=d\cdot v_i$。则
\begin{equation}
W_i=\bigl\{S_i+\xi u_i+\eta v_i:\ \xi\ge0,\ |\eta|\le\xi\tan\tau\bigr\}.
\tag{1}
\end{equation}
\end{lemma}

\begin{proof}
设 $d\neq0$。因 $\{u_i,v_i\}$ 为标准正交基，当 $\xi_i>0$ 时有
$\argvec d=\theta_i+\arctan(\eta_i/\xi_i)$。若 $\xi_i\le0$ 且 $d\neq0$，则
$\lvert\argvec d-\theta_i\rvert\ge90^\circ>\tau$，故 $d\notin W_i-S_i$；若 $\xi_i>0$，则
\[
d\in W_i-S_i\iff\bigl|\arctan(\eta_i/\xi_i)\bigr|\le\tau
\iff|\eta_i|\le\xi_i\tan\tau .
\]
又 $\xi=0\Rightarrow\eta=0$（否则仍有 $\lvert\argvec d-\theta_i\rvert=90^\circ$），故闭条件
$\xi\ge0,\ |\eta|\le\xi\tan\tau$ 恰在 $\xi=\eta=0$ 处补上基点 $S_i$。
\end{proof}

\begin{lemma}[楔形 = 两个半平面之交]\label{lem:halfplane}
令楔形两条边界的交点（楔顶）
\begin{equation}
A_i=S_i+\frac{1}{\cos\tau}u_i,
\end{equation}
则
\begin{equation}
W_i=\bigl\{P:\ n_i^{+}\!\cdot P\le n_i^{+}\!\cdot A_i\bigr\}
\cap\bigl\{P:\ n_i^{-}\!\cdot P\le n_i^{-}\!\cdot A_i\bigr\},
\qquad
n_i^{\pm}=\bigl(\sin(\theta_i\mp\tau),\ -\cos(\theta_i\mp\tau)\bigr).
\tag{2}
\end{equation}
\end{lemma}

\begin{proof}
由引理~\ref{lem:cone}，$W_i$ 的两条边界为射线
\[
S_i+t\,e_i^{\pm},\qquad e_i^{\pm}=\bigl(\cos(\theta_i\pm\tau),\ \sin(\theta_i\pm\tau)\bigr),\quad t\ge0,
\]
二者交于 $t=1/\cos\tau$ 处，即 $A_i$。向量 $n_i^{\pm}\perp e_i^{\pm}$ 且为单位向量；校验法向朝向：
轴向点 $S_i+u_i$ 满足
\[
n_i^{\pm}\!\cdot(S_i+u_i)-n_i^{\pm}\!\cdot A_i
=\cos\tau-\sec\tau=-\frac{\sin^2\tau}{\cos\tau}<0,
\]
而反向点 $S_i-u_i$ 给出 $-\cos\tau-\sec\tau>0$，故所取法向确指向楔形内部，两半平面的交即该角域。
\end{proof}

\begin{corollary}
$\mathcal R$ 是\textbf{闭凸集}，且是闭凸多边形（至多 $2n$ 条边）。
\end{corollary}

\section{有界性}

\begin{lemma}[逃逸方向]\label{lem:recession}
非零向量 $\xi$ 是 $\mathcal R_0=\bigcap_iW_i$ 的回收方向（即存在 $P_0\in\mathcal R_0$ 使
$P_0+t\xi\in\mathcal R_0,\ \forall t\ge0$）当且仅当
\begin{equation}
\bigl|\wrap(\theta_i-\argvec\xi)\bigr|\le\tau,\qquad i=1,\dots,n .
\tag{3}
\end{equation}
\end{lemma}

\begin{proof}
由式~(1)，$W_i$ 的回收锥为 $\{t\,e:\ |\argvec e-\theta_i|\le\tau,\ t\ge0\}$，而
$\bigcap_iW_i$ 的回收锥等于各回收锥之交。闭凸集有界当且仅当其回收锥为 $\{0\}$。
\end{proof}

\begin{theorem}[有界性判据]\label{thm:bounded}
记 $\Theta=\{\theta_1,\dots,\theta_n\}$ 的\textbf{最小包围弧}
\begin{equation}
\alpha(\Theta)=360^\circ-\max_{k}\bigl(\theta_{(k+1)}-\theta_{(k)}\bigr)_{360},
\end{equation}
其中 $\theta_{(k)}$ 为升序排列、下标循环、$(\cdot)_{360}$ 表示取 $[0,360^\circ)$ 内的环间隔。则
\begin{equation}
\boxed{\ \mathcal R_0\ \text{有界}\iff\alpha(\Theta)>2\tau\ }
\end{equation}
\end{theorem}

\begin{proof}
由引理~\ref{lem:recession}，$\mathcal R_0$ 无界 $\iff$ 存在方向 $\varphi$ 与所有 $\theta_i$ 的夹角均 $\le\tau$
$\iff$ 所有 $\theta_i$ 落入某条长度为 $2\tau$ 的弧内 $\iff\alpha(\Theta)\le2\tau$。取补即得。
\end{proof}

\begin{corollary}[两点情形]
$n=2$、$\Delta\theta=\lvert\wrap(\theta_2-\theta_1)\rvert<180^\circ$ 时 $\alpha=\Delta\theta$，故
\[
\mathcal R_0\ \text{有界}\iff\Delta\theta>2\tau=2^\circ .
\]
\end{corollary}

\begin{corollary}[模型适用性]\label{cor:boundedness}
当 $\Delta\theta\le2\tau$（两检测点几乎与干扰源共线）时 $\mathcal R_0$ 是无界双锥，直径无意义。
此时须叠加物理约束，例如
\[
\mathcal R=\mathcal R_0\cap\bigcap_{i}\{|P-S_i|\le R_{\max}\},
\quad 1000\le R_{\max}\le1500\ \text{m（附录~2(2)）}
\]
或 $|P|\le1800$ m（目标区域）。$\mathcal R$ 仍为紧凸多边形，下文结论不变。
\end{corollary}

\section{直径的精确计算}

\begin{definition}
$D=\diam\mathcal R=\sup_{P,Q\in\mathcal R}|P-Q|$。
\end{definition}

\begin{theorem}[直径 = 最大宽度]\label{thm:width}
对紧集 $K$，
\begin{equation}
D=\max_{\lVert\rho\rVert=1}w(\rho),\qquad
w(\rho)=\max_{P\in K}\rho\!\cdot P-\min_{P\in K}\rho\!\cdot P .
\tag{4}
\end{equation}
\end{theorem}

\begin{proof}
一方面，对任意 $P,Q\in K$ 与 $\lVert\rho\rVert=1$ 有 $|P-Q|\ge\rho\!\cdot(P-Q)$，故 $D\ge\max_\rho w(\rho)$。
另一方面，设 $P^*,Q^*$ 实现 $D$（紧集上可达），取 $\rho^*=(P^*-Q^*)/D$，则
$w(\rho^*)\ge\rho^*\!\cdot(P^*-Q^*)=D$，故 $\max_\rho w(\rho)\ge D$。
\end{proof}

\begin{theorem}[凸多边形情形的算法结论]\label{thm:vertex}
设 $\mathcal R$ 的顶点为 $V_1,\dots,V_m$（逆时针），则
\begin{equation}
\boxed{\ D=\max_{1\le i<j\le m}|V_i-V_j|\ }
\tag{5}
\end{equation}
\end{theorem}

\begin{proof}
固定 $Q$，函数 $P\mapsto|P-Q|$ 是凸函数；紧凸集上凸函数的最大值在极点（顶点）取得——否则若 $P$
位于某条边内部，则该边上 $|P-Q|$ 是凸的一元函数，其最大值必在端点。故 $P$ 可取为顶点；再对 $Q$
作同样论证即得式~(5)。
\end{proof}

\begin{remark}[算法步骤]
\textbf{(1)} 由式~(2) 得 $\mathcal R$ 的 $2n$ 条半平面 $n_j\!\cdot P\le c_j$，用半平面求交（逐步裁剪，
或按法向极角排序后单调队列）得顶点表 $V_1,\dots,V_m$，$m\le2n$，复杂度 $O(n^2)$。
\textbf{(2)} 按式~(5) 枚举顶点对求最大距离，$O(m^2)$；或先取凸包再用旋转卡壳 $O(m)$。
\textbf{(3)} 交叉校验：用式~(4) 在方向集 $\{\rho_{ij}=(V_j-V_i)/|V_j-V_i|\}$ 上取 $\max w(\rho)$，
与式~(5) 的结果必须一致（凸多边形上 $w$ 的局部极大只出现在某对顶点连线方向上）。
\end{remark}

\section{结论与反例：以直径为直径的圆能否覆盖定位区域}

\subsection{结论}

\boxed{\text{一般不能。}}

精确定量结论需引入最小覆盖圆半径
\begin{equation}
R_{\MEC}(K)=\min_{C}\max_{P\in K}|C-P| .
\end{equation}

\begin{theorem}\label{thm:mec}
设 $\diam K=D$，则
\begin{equation}
\frac D2\ \le\ R_{\MEC}(K)\ \le\ \frac{D}{\sqrt3},
\tag{6}
\end{equation}
且
\begin{equation}
\text{存在直径 }D\text{ 的圆覆盖 }K\iff 2R_{\MEC}(K)=D .
\tag{7}
\end{equation}
\end{theorem}

\begin{proof}
左端：$K$ 中有一对距离为 $D$ 的点，任何覆盖圆的半径必 $\ge D/2$。右端：由容斥（Jung）定理，
直径 $D$ 的集合可被半径 $D/\sqrt3$ 的圆覆盖（平面情形的最优常数）。
式~(7)：$R_{\MEC}\ge D/2$ 给出 $2R_{\MEC}\ge D$，故“直径 $D$ 的圆可覆盖”等价于 $2R_{\MEC}\le D$，
结合 $\ge$ 得 $=$。
\end{proof}

对内部非空的二维凸集，$2R_{\MEC}>D$ 通常严格成立，故\textbf{以定位区域直径 $D$ 为直径的圆不能覆盖
该区域}；可行的最小覆盖圆直径是
\[
2R_{\MEC}\in\Bigl[D,\ \tfrac{2}{\sqrt3}D\Bigr].
\]

\subsection{反例：等边三角形（凸集）}

取 $K$ 为边长 $D$ 的等边三角形。

\begin{enumerate}[label=\textbf{(\arabic*)}]
\item \textbf{直径}：三边均长 $D$，故 $\diam K=D$。
\item \textbf{中点圆不覆盖}：取实现直径的一对顶点 $A,B$，令 $M=(A+B)/2$，第三顶点为 $P$。
等边三角形中
\[
|MA|=|MB|=\frac D2,\qquad |MP|=\frac{\sqrt3}{2}D>\frac D2 ,
\]
故 $K\not\subset B\bigl(M,\tfrac D2\bigr)$。
\item \textbf{任何同直径圆都不覆盖}：等边三角形的外接圆即其最小覆盖圆，半径 $\dfrac{D}{\sqrt3}$，于是
\[
2R_{\MEC}=\frac{2}{\sqrt3}D\approx1.1547\,D>D ,
\]
由式~(7) 得否定结论。
\item \textbf{紧性}：$\dfrac{2R_{\MEC}}{D}=\dfrac{2}{\sqrt3}$ 恰为常数上界（等边三角形取到）；
而细长平行四边形可使该比值任意接近 $1$，故 $D$ 与 $2R_{\MEC}$ 之间无更好的普适关系。
\end{enumerate}

\begin{remark}
等边三角形是凸集，与“定位区域为凸多边形”的语境完全相容，故该反例有效。
\end{remark}

\subsection{模型层面的注记}

若题目本意是“\textbf{给定圆心}（如取区域形心或直径端点中点）、以区域直径为直径”，则结论同样为否定。
事实上，对任意点 $M$，取实现直径的 $A,B$ 与任一点 $P\in K$，由
$|MA|=|MB|=D/2$ 且 $|AP|,|BP|\le D$，利用平行四边形恒等式
\[
|MP|^2=\frac{|AP|^2+|BP|^2}{2}-\frac{|AB|^2}{4}\le\frac{D^2+D^2}{2}-\frac{D^2}{4}=\frac34D^2,
\]
得 $|MP|\le\dfrac{\sqrt3}{2}D$，该上界仍严格大于 $D/2$。故任何“圆心约定 $+$ 直径为 $D$”的方案都不能
保证覆盖；\textbf{正确做法是报告最小覆盖圆直径 $2R_{\MEC}$（或以其作为搜索域尺度）}。

\section{推导链与结论一览}

\begin{table}[h]
\centering
\small
\renewcommand{\arraystretch}{1.3}
\begin{tabular}{>{\bfseries}lll}
\toprule
步骤 & 结论 & 依据 \\
\midrule
1. 建模 & $\mathcal R=\bigcap_i\{\lvert\wrap(\argvec(P-S_i)-\theta_i)\rvert\le\tau\}$，$G\in\mathcal R$ & 定义、命题~1 \\
2. 线性化 & $\mathcal R=\bigcap_i\{n_i^{\pm}\!\cdot P\le n_i^{\pm}\!\cdot A_i\}$：闭凸多边形 & 引理~1、2，推论~1 \\
3. 可解性 & 有界 $\iff\alpha(\Theta)>2\tau$；$n=2$ 时 $\iff\Delta\theta>2^\circ$ & 引理~3、定理~\ref{thm:bounded}、推论 \\
4. 直径 & $D=\max_\rho w(\rho)=\max_{i<j}|V_i-V_j|$，$O(n^2)$（或 $O(n)$ 卡壳） & 定理~\ref{thm:width}、\ref{thm:vertex} \\
5. 第二问 & \textbf{不能覆盖}；$D\le2R_{\MEC}\le\frac{2}{\sqrt3}D$，反例为等边三角形 & 定理~\ref{thm:mec} + 反例 \\
\bottomrule
\end{tabular}
\end{table}

\end{document}
